\documentclass[fontset=none,11pt,a4paper]{ctexart}
\setCJKmainfont{AR PL SungtiL GB}[BoldFont=AR PL UMing CN,AutoFakeBold=2.5]
\setCJKsansfont{AR PL UMing CN}
\setCJKmonofont{AR PL UMing CN}
\usepackage{geometry}
\usepackage{booktabs}
\usepackage{multirow}
\usepackage{hyperref}
\usepackage{enumitem}
\usepackage{xcolor}
\geometry{margin=2.5cm}
\title{多模态Agent工具调用Post-Training：\\问题、方法与系统总结}
\author{}
\date{\today}

\begin{document}
\maketitle

\begin{abstract}
本项目以Qwen2.5-VL-3B/7B为基座模型，通过SFT/DPO LoRA微调，训练VLM在发票核验、图表计算、图文一致性等场景下严格遵循\texttt{vision\_parse $\to$ retrieve\_docs/python\_exec $\to$ final\_answer}的工具调用流程。本文系统梳理了项目全过程中遇到的问题、采取的应对策略、取得的效果及从中提炼的关键教训。核心发现：(1) 单阶段SFT在同时优化多目标（格式纪律、工具选择、答案质量、拒答边界）时存在根本性局限；(2) 更大规模模型（7B vs 3B）虽然文档理解能力更强，但对训练数据改动的敏感度呈数量级增加；(3) 当前最优模型为V4(3B)和V7B-clean(7B)，分别在工具纪律和文档理解上达到天花板。
\end{abstract}

\section{项目背景}

\subsection{任务定义}
给定一张图片和一个自然语言问题，VLM需要按以下流程完成任务：
\begin{enumerate}[label=(\arabic*)]
    \item 判断任务类型（发票核验/图表计算/图文一致性/品牌识别/简单描述）
    \item 如需视觉信息，调用 \texttt{vision\_parse(mode)} 进行OCR/图表读数/图像描述
    \item 如需规则依据，调用 \texttt{retrieve\_docs(query)} 检索规则文档
    \item 如需数学计算，调用 \texttt{python\_exec(code)} 完成计算
    \item 输出包含依据的 \texttt{<final\_answer>} 或合理拒答
\end{enumerate}

每轮对话模型只能输出一个 \texttt{<tool\_call>} 或一个 \texttt{<final\_answer>}，不得同时输出多个动作或裸文本。

\subsection{场景与数据集}
\begin{itemize}[nosep]
    \item \textbf{训练数据}：50条种子数据（人工构造mock vision\_parse结果） + ChartQA(500条) + DocVQA(200条) + CORD(203条) = 1133条（v4配置）
    \item \textbf{评测数据}：50条，分5类------document\_validation(15)、image\_text\_consistency(12)、chart\_calculation(10)、uncertainty\_refusal(8)、direct\_vqa(5)
    \item \textbf{工具}：vision\_parse（真实VLM推理，mock回退）、retrieve\_docs（关键词匹配3份规则文档）、python\_exec（安全正则白名单eval）
    \item \textbf{GPU}：NVIDIA RTX 4080 SUPER 48GB
\end{itemize}

\section{实验全景}

\subsection{模型版本一览}

\begin{table}[htbp]
\centering
\begin{tabular}{lllllp{3cm}}
\toprule
版本 & 基座 & 数据量 & LoRA & 关键改动 & 结果 \\
\midrule
V2 & 3B & 446 & r=16 & mock数据 + ChartQA & 基准线 \\
V3 & 3B & 1270 & r=8 & +DocVQA+CORD+recovery & 过度工具调用 \\
V4 & 3B & 1133 & r=16 & \textbf{真实VLM vision\_parse} & \textbf{3B最优} \\
V5 & 3B & 1611 & r=16 & +800条增强数据 & kw持平(82\%) \\
V6 & 3B & 1611 & r=16 & 修refusal模板 & kw退步(70\%) \\
\addlinespace
DPO v1 & 3B & 58对 & -- & 在V2上DPO & 格式崩溃(60\%) \\
DPO v4-mt & 3B & 13对 & -- & 多轮DPO & 格式崩溃(14\%) \\
SFT微调 & 3B & 13条 & r=8 & V4上微调 & 格式崩溃(22\%) \\
\addlinespace
V7 & 7B & 1611 & r=16 & 增强数据 & 崩溃(14\% tool) \\
V7B-clean & 7B & 1133 & r=32 & \textbf{纯V4数据} & \textbf{7B最优} \\
V7B-opt & 7B & 1133 & r=64 & 高rank & 崩溃(16\% tool) \\
V7B-v2 & 7B & 1145 & r=32 & +9条refusal & 崩溃(16\% tool) \\
\addlinespace
V8 & 3B & 1186 & r=16 & +53条工具决策样本 & kw退步(50\%) \\
\bottomrule
\end{tabular}
\caption{所有模型版本及其关键改动与结果。绿色=有效，红色=失败/退化，黄色=持平。}
\label{tab:models}
\end{table}

\subsection{核心评测指标}

采用模糊数值匹配后的最终指标（50条评测集）：

\begin{table}[htbp]
\centering
\begin{tabular}{lcccccl}
\toprule
模型 & Tool Hit & Keyword & Refusal & A/B/C/D & 无enf变化 \\
\midrule
V2 (3B, mock) & 50/50 & 41/50 & 48/50 & 45/0/0/5 & -- \\
V4 (3B, real) & 50/50 & 38/50 & 47/50 & 45/1/0/4 & tool仍100\% \\
V7B (7B) & 46/50 & 43/50 & 46/50 & 41/0/4/5 & tool升至49/50 \\
\bottomrule
\end{tabular}
\caption{最佳模型详细对比。A=需要工具且正确调用，B=不需工具且未调用，C=需要但漏调，D=不需但多调。}
\label{tab:results}
\end{table}

\section{问题与应对}

\subsection{问题1：训练-推理不匹配（mock vs 真实VLM）}

\textbf{现象}：训练数据中vision\_parse结果是人工构造的（如\texttt{OCR结果：Invoice No=INV-1001}），而外部数据集（ChartQA/DocVQA/CORD）甚至使用占位文本（\texttt{图表已解析；请结合图表视觉证据...}）。模型从未见过真实VLM输出分布。

\textbf{应对}：
\begin{itemize}[nosep]
    \item 新建\texttt{scripts/rebuild\_seed\_with\_real\_vlm.py}，加载真实Qwen2.5-VL-3B模型，对46张种子图片逐张做vision\_parse推理，替换mock结果
    \item 同样处理751张外部数据集图片（ChartQA 500 + DocVQA 200 + CORD 51），对CORD的152次重复推理使用image+mode级缓存
\end{itemize}

\textbf{效果}：797张图片全部完成真实VLM处理，0错误。V4模型从mock数据迁移到真实数据后，tool hit保持100\%，document\_validation从10/15提升至15/15（7B）。

\subsection{问题2：大图导致OOM}

\textbf{现象}：DocVQA图片最大1.8MB，CORD图片最大7.4MB（vs 种子图片~10KB）。Qwen-VL动态分辨率机制对大图生成海量visual token，32GB显存OOM。尝试bf16→fp16→fp16+gradient\_checkpointing→max\_length=1024→PYTORCH\_CUDA\_ALLOC\_CONF全部失败。

\textbf{应对}：将所有外部图片resize到最大1024px（最长边），243/751张图片被压缩。训练数据JSONL路径更新至\texttt{images\_resized/}目录。

\textbf{效果}：bf16 + 3B + 1133行顺利训练，29GB/32GB。7B + gradient\_checkpointing使用20GB/48GB。48GB升级后仍无法用原图训练7B（7B+r=64+原图OOM at 41.6GB），resize是必须的。

\subsection{问题3：v3过度工具调用}

\textbf{现象}：v3在不需要工具的场景（如简单商品描述）也强行调用工具。根因是recovery变体教模型"宁可多调，不能漏调"；ChartQA/CORD上采样5x过度放大第二工具信号。

\textbf{应对}：
\begin{itemize}[nosep]
    \item recovery\_repeat: 2→1
    \item python\_exec\_repeat: 5→2
    \item retrieve\_docs\_repeat: 5→2
    \item recovery提示词与runtime.py的\_missing\_required\_tool\_prompt对齐
\end{itemize}

\textbf{效果}：v4工具命中100\%，同时D类（不需但多调）从v3的显著降低。

\subsection{问题4：DPO反复崩溃}

\textbf{现象}：三次DPO尝试全部导致格式纪律崩溃------模型开始输出裸文本、一次输出多个动作、或输出空字符串。
\begin{itemize}[nosep]
    \item DPO单轮（64对）：tool hit 100\%→60\%
    \item DPO多轮（13对）：tool hit→14\%，输出空字符串
    \item SFT微调（13条）：tool hit→22\%
\end{itemize}

\textbf{根因分析}：
\begin{enumerate}[label=(\arabic*),nosep]
    \item \textbf{训练-推理结构不匹配}：DPO数据是单轮（user→final\_answer），但推理是多轮（user→tool\_call→tool\_result→user\_continue→...→final\_answer）。单轮训练冲垮了SFT建立的多轮格式纪律。
    \item \textbf{数据量不足}：13或64条偏好对远不够覆盖多轮工具调用场景的变体。
    \item \textbf{格式无保护}：DPO loss只优化内容偏好，不惩罚格式错误。模型可以在降低DPO loss的同时完全抛弃格式约束。
    \item \textbf{小数据过拟合}：13条×3 epochs的梯度信号没有足够多样性维持格式纪律。
\end{enumerate}

\textbf{应对}：放弃DPO，回到SFT做数据驱动优化。

\textbf{关键教训}：在严格格式约束的任务上，DPO必须满足三个条件------(1) 500+条多轮偏好对、(2) SFT+DPO混合loss保留格式、(3) 偏好对必须是完整多轮轨迹。当前数据量不满足这些条件。

\subsection{问题5：7B对训练数据改动的极端敏感性}

\textbf{现象}：v7b-clean（纯V4数据，r=32）成功训练出92\% tool hit。但任何对训练数据的改动------加9条refusal样本(V7B-v2)、加818条增强样本(V7)、提高rank到64(V7B-opt)------全部导致模型崩溃到14-16\% tool hit。

\textbf{诊断}：
\begin{itemize}[nosep]
    \item 失败的模型输出模式完全一致：直接输出final\_answer而不调用任何工具
    \item 答案内容通常是mock数据中的数值（如"金额是107.50"），而非真实VLM结果
    \item 崩溃阈值：仅需改动<1\%的训练数据（9/1133=0.8\%）
\end{itemize}

\textbf{可能原因}：
\begin{enumerate}[label=(\arabic*),nosep]
    \item 7B更强的记忆能力导致更倾向背答案而非学流程
    \item 训练数据的工具调用样本格式高度统一，任何"不同格式"的样本（如直接回答、不同推理链）都会破坏模型对这一格式的专注
    \item LoRA adapter容量有限，新样本的梯度信号与原有格式信号竞争
\end{enumerate}

\textbf{应对}：保留v7b-clean作为7B最优模型，放弃在7B上增量添加训练数据的策略。

\textbf{关键教训}：7B对该训练数据配置存在一个极窄的收敛盆地。一旦训练数据分布发生任何偏移，模型就会跌出这个盆地进入"裸回答"模式。这一现象本身是值得研究的模型行为。

\subsection{问题6：单阶段SFT的多目标冲突}

\textbf{现象}：所有"修补"特定问题的尝试都导致其他能力退化：
\begin{itemize}[nosep]
    \item V8加47条直接回答样本→tool仍100\%但kw从38跌至25，D类不降反增
    \item V6修refusal模板→kw从82\%跌至70\%
    \item V5加800条增强数据→kw持平但不提升
\end{itemize}

\textbf{根因}：单阶段SFT需要同时优化四个互相矛盾的目标：
\begin{enumerate}[label=(\arabic*),nosep]
    \item 严格格式纪律（每轮只输出一个动作）
    \item 正确的工具选择（调用什么、何时调用）
    \item 准确的答案内容（数值精度、关键词）
    \item 恰当的拒答边界（何时该拒答）
\end{enumerate}
这些目标在梯度空间中是竞争关系------优化"直接回答"会削弱"工具调用"，优化"内容准确"可能以"格式正确"为代价。

\subsection{问题7：Eval标签与实际模型能力不匹配}

\textbf{现象}：
\begin{itemize}[nosep]
    \item 5个document\_validation样本eval标注为"缺少Tax ID"，但真实VLM从图中读到了Tax ID。V4/V7B的"字段完整"答案实际正确，但被计为keyword miss。
    \item 3个可读logo（ACME/LUNA/VOLT）eval要求先调vision\_parse，但7B已学会直接识别，导致虚假tool miss。
    \item 图表计算的eval要求精确数值匹配"25.00"，模型输出"25.0\%"被判定为keyword miss。
\end{itemize}

\textbf{应对}：修正eval标签（5个invoice改为"完整"、3个可读logo放宽gold\_tools、图表使用模糊数值匹配）。

\textbf{效果}：V4 keyword从72\%→76\%（修正eval后），V7B从46/50 tool→49/50 tool。

\textbf{关键教训}：从mock数据切换到真实VLM后，eval标签必须同步更新。基于mock预期的gold label在真实VLM场景下不再准确。

\subsection{问题8：enforcement的副作用}

\textbf{现象}：V7B在有enforcement时tool hit=46/50，关掉enforcement后反而升至49/50。Enforcement在阻碍7B------当7B正确识别可读logo并直接回答时，enforcement强制它调vision\_parse，导致max\_steps截断。

\textbf{应对}：对可读logo放宽eval的gold\_tools要求，对不可读logo保持严格要求。

\textbf{关键教训}：enforcement是训练阶段的拐杖，不是推理阶段的必需品。当模型能力超过enforcement规则时，enforcement会成为负优化。

\section{关键发现}

\begin{enumerate}[label=\textbf{\arabic*}. , leftmargin=*]

\item \textbf{训练数据分布存在脆弱平衡}：V4训练数据的成功源自其精确的工具调用-答案分布。任何扰动------哪怕是0.8\%的数据改动------都会打破这个平衡，导致模型性能崩溃。这一现象在7B上尤其显著。

\item \textbf{模型规模与数据敏感度正相关}：3B可以容忍800条额外增强数据（V5持平），7B连9条refusal样本都承受不住。更大模型对训练数据分布偏移更敏感，这是VLM post-training中的重要建模约束。

\item \textbf{DPO不适用于严格格式任务（小数据场景）}：在每轮只能输出单一动作的格式约束下，DPO的last-turn log-prob优化与多轮tool-calling的推理结构不兼容。需要大规模多轮偏好对和混合loss才能安全使用。

\item \textbf{Enforcement是双刃剑}：train-time enforcement帮助模型建立工具调用基线，但inference-time enforcement可能阻碍已经学会正确判断的模型。

\item \textbf{3B和7B存在互补的能力分布}：3B工具纪律更强（100\% vs 92\%），7B文档理解更强（doc 15/15 vs 10/15, img 12/12 vs 9/12）。两个规模有各自的能力天花板，没有单一的"全优"配置。

\item \textbf{LoRA rank存在最优区间}：3B在r=16最优，7B在r=32最优。rank太低（r=8）表达能力不足，rank太高（r=64）导致过拟合和模式坍缩。

\end{enumerate}

\section{方法论反思}

\subsection{为什么大部分优化尝试失败了}

回顾从V4到V9的迭代，一个反复出现的模式是：在已经收敛到良好状态的模型上做增量修改→模型崩溃→回到原点。这指向一个方法论问题：

\begin{center}
\textbf{单阶段SFT在严格格式+多目标任务上存在根本性天花板。}
\end{center}

单阶段SFT将四个目标（格式、工具选择、答案质量、拒答边界）压缩进同一个loss函数。模型到达的"最优点"实际上是一个平衡点，而非每个目标上的最优。任何针对单一目标的"修补"都会打破这个平衡。

\textbf{解决路径}：需要将多目标解耦------
\begin{itemize}[nosep]
    \item 两阶段训练：第一阶段纯格式，第二阶段加内容和判断
    \item 或RL（GRPO）：将格式硬编码为reward函数中的硬约束，让模型在格式安全区内优化内容和判断
\end{itemize}
其中RL更符合本任务的特点（有明确的、可量化的reward信号），且格式约束可以天然地写进reward函数。

\subsection{当前最佳模型}

\begin{table}[htbp]
\centering
\begin{tabular}{lcccl}
\toprule
 & V4 (3B) & V7B-clean (7B) & 互补？ \\
\midrule
Tool Hit & \textbf{50/50} & 46/50 (49/50 w/o enf) & V4更优 \\
Keyword & 38/50 & \textbf{43/50} & V7B更优 \\
Document Validation & 10/15 & \textbf{15/15} & V7B显著优 \\
Image Consistency & 9/12 & \textbf{12/12} & V7B显著优 \\
Chart Calculation & 8/10 & 7/10 & 持平 \\
Refusal & \textbf{47/50} & 46/50 & V4略优 \\
数据敏感性 & \textbf{低} & 极高 & V4更鲁棒 \\
\bottomrule
\end{tabular}
\caption{V4与V7B-clean全面对比。}
\label{tab:comparison}
\end{table}

两个模型各有所长，且能力高度互补。V4适合需要100\%工具纪律的场景，V7B适合文档理解密集的场景。

\section{后续方向}

\begin{enumerate}[label=\textbf{\arabic*}. , leftmargin=*]

\item \textbf{RL (GRPO) 优化}：用V4或V7B-clean作为base model，定义包含格式、工具选择、答案质量的多维reward函数，通过GRPO在保持格式约束的同时优化内容和判断能力。这是突破当前SFT天花板的最有希望的路径。

\item \textbf{模型规模的工具学习动力学}：系统对比3B/7B在相同数据下的工具调用行为差异------为什么7B文档理解更强但工具纪律更弱？为什么7B对数据扰动更敏感？这些问题的答案对理解VLM的工具学习机制有基础性贡献。

\item \textbf{数据污染阈值研究}：在纯净训练数据中逐步混入不同比例的"直接回答"样本（0.1\%/0.5\%/1\%/5\%），测量每种配比下的工具命中率变化，量化得到"格式崩溃阈值"。

\item \textbf{无enforcement的工具调用自主性}：当前eval仍以enforcement模式为主。真正的研究目标应该是让模型无需runtime兜底就能自主完成正确的工具调用------这是post-training研究的本质问题。

\item \textbf{跨模型蒸馏}：V4有完美的工具纪律，V7B有最强的文档理解。将V4的工具调用能力蒸馏到7B上，可能融合两者的优势。

\end{enumerate}

\section{结论}

本项目围绕"VLM工具调用post-training"进行了系统性的探索，训练并评测了十余版模型。核心发现是：

\begin{enumerate}[nosep]
    \item 真实工具接入（V4的vision\_parse用真实VLM替代mock）是唯一一次质的飞跃；
    \item 此后所有增量优化均在单阶段SFT框架内失败------数据扰动导致格式崩溃、多目标冲突导致顾此失彼；
    \item 7B在文档理解上显著优于3B，但对训练数据变化的敏感度极高（0.8\%的改动即可导致崩溃）；
    \item DPO在严格格式的小数据场景下反复失败，需要大规模多轮偏好对和混合loss；
    \item 当前最优模型V4(3B)和V7B-clean(7B)分别在工具纪律和文档理解上达到天花板；
    \item 下一步最值得尝试的方向是RL(GRPO)------用显式reward函数解耦格式、工具选择和答案质量的优化。
\end{enumerate}

\end{document}
